% Copyright (c) 2024 by Daniel R. Driver.
% !TEX encoding = UTF-8 Unicode
% !TEX TS-program = XeLaTeX

\documentclass[titlepage]{article}

% This document presumes a file structure and set of inputs that are
% available at: git@github.com:danieldriver/syllabi.git

\newcommand\policy{../policy}
\newcommand\incl{../includes}
\ProvidesFile{variables.tex}[2022/09/11 v2.3 -- Syllabus variables]

\usepackage{xspace} % make manual spaces (like \mycmd\ ) unnecessary
\usepackage{xifthen} % provides \isempty test

% variables for internal use
\newcommand\prof{}
\newcommand\pdegree{}
\newcommand\pphone{}
\newcommand\pemail{}
\newcommand\poffice{}
\newcommand\phours{}
%
\newcommand\ccode{}
\newcommand\ctitle{}
\newcommand\cseries{}
\newcommand\cversion{}
\newcommand\csemester{}
\newcommand\cmeetson{}
\newcommand\cmeetsat{}
\newcommand\cmeetsin{}
\newcommand\cwebsite{}
\newcommand\cdescrip{}
\newcommand\cprereqs{}
\newcommand\edobject{}

% in case of fully online, BTh, or undergraduate courses - https://tex.stackexchange.com/a/5896
\newif\ifonline
\newcommand\Int[2]{\ifonline#1\else#2\fi}
\newif\ifbth
\newcommand\bth[2]{\ifbth#1\else#2\fi}
\newif\ifUG
\newcommand\UG[2]{\ifUG#1\else#2\fi}

% commands for setting variables in the preamble
\newcommand\professor[2][PhD]{
  \renewcommand\pdegree{#1\xspace}
  \renewcommand\prof{#2\xspace}}
\newcommand\phone[1]{
  \renewcommand\pphone{\addfontfeatures{Numbers=Monospaced}#1\xspace}}
\newcommand\email[1]{
  \renewcommand\pemail{\href{mailto:#1}{#1}\xspace}}
\newcommand\officehours[2][AST Library/Teams]{
  \renewcommand\poffice{#1\xspace}
  \renewcommand\phours{\href{https://outlook.office.com/bookwithme/user/b97b147a2e5c43b4a0e200c2cb630b77@astheology.ns.ca?anonymous}{#2}\xspace}%
  }
%
\newcommand\coursecode[2][1.0]{
  \renewcommand\cversion{#1\xspace}
  \renewcommand\ccode{#2\xspace}}
\newcommand\coursetitle[2][]{
  \ifthenelse{\isempty{#1}}%
    {}% do nothing if #1 is empty, else:
    {\renewcommand\cseries{#1\\[1ex]}}
  \renewcommand\ctitle{#2\xspace}}
\newcommand\semester[1]{
  \renewcommand\csemester{#1\xspace}}
\newcommand\meets[3]{
  \newcommand\AM{\textsc{am}}
  \newcommand\PM{\textsc{pm}}
  \renewcommand\cmeetson{#1\xspace}
  \renewcommand\cmeetsat{\Int{From 9:00 \AM}{#2}\xspace}
  \renewcommand\cmeetsin{\Int{\href{https://smu.brightspace.com/d2l/login}{Brightspace}}{#3}\xspace}}
\newcommand\website[2][]{
  \renewcommand\cwebsite{
  \ifthenelse{\isempty{#1}}%
    {\href{https://#2}}% copy url if #1 is empty, else
    {\href{#1}}% use specified url and print #2
  {#2}\xspace}}
\newcommand\cdescription[2][BF 1001; RM 1000 or GTRS 6000]{
  \renewcommand\cprereqs{#1}
  \renewcommand\cdescrip{#2\par}}
\newcommand\objectives[1]{
  \renewcommand\edobject{#1\par}}


%\onlinetrue % \Int{true}{false}
\coursecode[3.7]{HB 2201}
\coursetitle[Hebrew Bible]{Beginning Biblical Hebrew I}

% Taught as HEBR 2013 in:
%   - Fall 2009
%   - Fall 2011
%   - Fall 2013
%   - Fall 2015
% Taught as HB 2201 in:
%   - Fall 2018
%   - Summer 2020 (on Teams)
%   - Summer 2022 (on campus)
%   - Fall 2024 (hybrid)

\professor{Daniel R. Driver}
\phone{902-425-7051}
\email{daniel.driver@astheology.ns.ca}
\officehours{Mondays, 1:00--3:00 \PM}

\semester{Fall 2024}
\meets{Thursdays}% \meets{on}{at}{in}
      {2:00--4:30 \PM}% if online, it prints "From 9:00 AM"
      {Classroom 3}% if online, it prints "Brightspace" with link
\website{danieldriver.com}
\cdescription[none. The course is required for HB 2202 and advanced biblical Hebrew]{% copy from the current Academic Calendar; [] for prereqs

	For students who wish to read, hear, and even (to an extent) produce
	Hebrew, this course offers a textually immersive introduction to
	classical Hebrew. Grammatical features like morphology, phonology,
	and syntax are learned inductively, through a progression of
	readings drawn from Genesis. Techniques of second-language
	acquisition are also used to help develop language competency.
	Continuation in HB 2202 is strongly recommended. This course is open
	to for-credit students and to Continuing Education participants.

}% end of course description
\objectives{% recall Bloom's taxonomy: http://www.celt.iastate.edu/teaching/RevisedBlooms1.html
	The basic goal of this course is to begin to learn the rudiments of
	biblical Hebrew. The course is designed as the first part of a
	two-semester series. Since some may not continue in the Fall,
	however, \ccode is being offered with two different options for
	outcomes.

	Option 1: Learn the alphabet, vowels, and syllables with the aim of
	continuing to learn Hebrew in HB 2202. Students will become
	familiar with the components that can be added to the front and back
	of Hebrew nouns and adjectives so that they can gain a basic
	understanding of the relationship between parsings in electronic
	sources (Accordance, Logos) and parallel translations. The verb
	system will be introduced briefly with the same purpose. On this
	path, students will read and translate a few passages from Genesis
	so that they can begin to read the Hebrew Bible with understanding.
	Students will learn about the limits of what they know so that they
	are less likely to misapply what they learn in church contexts. This
	option should be the default for students taking the course for
	credit.

	Option 2: Learn the alphabet, vowels, and syllables with the aim of
	continuing to learn Hebrew on your own. Students will hear and read
	the same overviews as in Option 1, but may choose to work through
	the assigned exercises at a slower pace. The goal is to gain enough
	experience to become self-learners who can complete assigned
	readings and translations as they are able. Students will build
	enough confidence and momentum to progress through the grammar and
	workbook on their own schedule. This option caters to the objectives
	of some continuing education students, for example, but without
	imposing limits. Note that, in a past iteration of the course, the
	highest achiever overall was an auditor.

	A goal in every case is to whet the appetite for learning Biblical
	Hebrew for long-term study and use. Language is communal. In the
	near term, therefore, every student will also have the opportunity
	to become part of a community of language learners.
}% end of learning objectives

\ProvidesFile{preamble.tex}[2014/09/02 v1.1 -- Syllabus preamble]

% basic typography
\usepackage{fontspec,realscripts}
\defaultfontfeatures{Ligatures=TeX}
\setmainfont[%
  Numbers={Proportional,OldStyle},
  SmallCapsFeatures={RawFeature=+c2sc;+smcp,LetterSpace=1}]{Scala OT}
\newfontfamily\lining[Numbers={Lining}]{Scala OT}
\setsansfont[%
  SmallCapsFeatures={RawFeature=+c2sc;+smcp,LetterSpace=1}]{Scala Sans Pro}
\setmonofont[Scale=MatchLowercase]{Menlo}
\usepackage{sectsty}
\allsectionsfont{\sffamily\scshape}
\frenchspacing
\setlength{\emergencystretch}{3em} % prevent overfull lines

% custom font size and leading
\renewcommand\tiny{\fontsize{6}{9}\selectfont}
\renewcommand\scriptsize{\fontsize{7}{10}\selectfont}
\renewcommand\footnotesize{\fontsize{8}{10}\selectfont}
\renewcommand\small{\fontsize{9}{11}\selectfont}
\renewcommand\normalsize{\fontsize{10}{12}\selectfont}% base size
\renewcommand\large{\fontsize{11}{13}\selectfont}
\renewcommand\Large{\fontsize{12}{14}\selectfont}
\renewcommand\LARGE{\fontsize{16}{18}\selectfont}% "course syllabus \\ semester" benefits from more lead
\renewcommand\huge{\fontsize{20}{22}\selectfont}
\renewcommand\Huge{\fontsize{24}{26}\selectfont}

% layout packages: page, logo, tables
\usepackage[scale={0.6,0.8},
            xetex]{geometry}
\usepackage{graphicx}
\usepackage{array}     % allow insertions of column styling with >{}
\usepackage{booktabs}  % elegant horizontal rules in tables

% custom macros for a session count in the schedule of readings
\newcounter{session}
\newcounter{columns}
\newcounter{spancolumns}
\newcounter{midcolumns}
\newcounter{courseunit}
\newcommand\setcolumncount[2][0]{ % optionally set count to other than 0,
  \setcounter{session}{#1}        % e.g. to -1, or to a standing count
  \setcounter{columns}{#2}
  \setcounter{spancolumns}{#2-1}
  \setcounter{midcolumns}{#2-2}}
\newcommand\sessioncount{\stepcounter{session}\arabic{session}}
\newcommand\sessionskip[1]{\multicolumn{1}{@{}r@{ }}{#1}}
\newcommand\unit[1]{\multicolumn{\thecolumns}{c}{%
  \scshape\stepcounter{courseunit}\roman{courseunit}. #1}}
\newcommand\noclass[1]{\sessionskip{---.}&\multicolumn{\themidcolumns}{@{}l}{\itshape No Class: #1}}
\newcommand\rarr{\Large\char"EA31\hspace*{0.5pt}} % pointer finger w/ reminder tie in Scala OT
\newcommand\reminder[2]{\sessionskip{\rarr}&\multicolumn{\thespancolumns}{@{}l}{#1 \hfill #2}}

% color to match Tyndale's branding
\usepackage[usenames]{xcolor}
% predefined: black, white, red, green, blue, cyan, magenta, yellow
\definecolor{TyndaleURLs}{HTML}{0062A0} % links on tyndale.ca
\definecolor{TyndaleBlue}{cmyk}{1,1,0,.32}
\definecolor{TyndaleGold}{cmyk}{0,.27,1,0}
\definecolor{TyndaleRed}{cmyk}{0,1,.99,.04}
\definecolor{TyndaleBlack}{cmyk}{0,0,0,1}
\definecolor{TyndaleGreen}{cmyk}{.45,0,1,.24}
\definecolor{TyndaleOrange}{cmyk}{0,.79,1,0}
\definecolor{TyndaleAqua}{cmyk}{.47,0,.24,0}
\definecolor{TyndaleYellow}{cmyk}{.03,.03,.35,0}

% metadata (assumes a host of definitions are made in the main file)
\usepackage[setpagesize=false,     % leave this to geometry
            hyperfootnotes=false,  % fragile and distracting
            xetex]{hyperref}
\hypersetup{breaklinks=true,       % allow link text to break across lines
            colorlinks=true,       % colorlinks resets pdfborder to 0 0 0
            urlcolor=TyndaleURLs,  % for external links
            linkcolor=TyndaleRed,  % for normal internal links
            citecolor=TyndaleGold, % for bibliographical citations in text
            pdfauthor={\prof},
            pdftitle={\ccode: \ctitle},
            pdfsubject={Tyndale UC, \csemester},
            pdfcreator={github.com/danieldriver/syllabus}}
\urlstyle{same}                    % don't use monospace font for urls

% custom footlines
\usepackage{fancyhdr}
\pagestyle{fancy} % turn it on
\fancyhf{}        % reset everything
\renewcommand{\headrulewidth}{0pt} % remove header line as well
\lfoot{\sffamily\scshape\footnotesize\ctitle, \csemester}
\rfoot{\sffamily\scshape\footnotesize\prof\quad\thepage}

% language support: polyglossia needs to go last
\usepackage{polyglossia}
\setdefaultlanguage{english}
\setotherlanguage{hebrew}
\newfontfamily\hebrewfont[Script=Hebrew]{SBL BibLit}

% protect positioning of margin table in policy/grades -- needs to come after polyglossia
\usepackage{marginfix}

% gratuitous with custom title page, but useful as a fallback
\title{\ccode: \ctitle}
\author{\professor}
\date{\semester}


\begin{document}
\ProvidesFile{title.tex}[2013/09/06 v1.0 -- Syllabus title page]

\begin{titlepage}
  \begin{center}

    \begin{minipage}{\textwidth}
      \parbox[t]{0.5\textwidth}{
        \mbox{}\\[-\baselineskip] % dummy line to align parboxes
        \includegraphics[width=0.5\textwidth]{.syllabus/includes/TyndaleUC}}
      \hfill
      \parbox[t]{0.4\textwidth}{
        \raggedleft\LARGE\sffamily Course Syllabus\\
        \semester}
    \end{minipage}

    \vfill

    {\sffamily\textsc{\LARGE\MakeLowercase\coursecode}\\[2ex]
      \bfseries\itshape\Huge\coursetitle}

    \vfill

    \begin{tabular}{>{\bfseries}rl>{\bfseries}rl}
      \toprule
      Instructor & \professor                   & Course  & \liningnums{\coursecode}~\texttt{(v\version)} \\
      \midrule
      Phone      & \texttt{\phone}              & Meets   & \meeton                                       \\
      Email      & \href{mailto:\email}{\email} & Time    & \meetat                                       \\
      Office     & \officenumber                & Room    & \meetin                                       \\
      Hours      & \officehours                 & Website & \href{http://\website}{\website}              \\
      \bottomrule
    \end{tabular}

    \vfill

    \begin{description}
      \item[Commuter Hotline]
        Class cancellations due to inclement weather or illness will
        be announced on the commuter hotline at \texttt{416.226.6620
        x2187}. Alternately, weather cancellation information is posted
        at \href{http://tyndale.ca/weather}{tyndale.ca/weather}.
      \item[MyTyndale.ca]
        This course may have materials stored on its website, such as
        handouts or readings that may be needed in order to complete
        assignments. Students are responsible for checking these course
        pages on a regular basis. Here, too, students are able to view
        their grades throughout the semester. For more information see
        Section~\ref{mytyndale}, below.
      \item[Mail]
        Students are responsible for information communicated through
        their campus mailboxes and student e-mail accounts. A mailbox
        directory hangs beside the mailboxes. For more information
        contact the Registrar's office.
    \end{description}

  \end{center}

  \section{Course Description}
  \label{description}

  \emph{From the Academic Calendar:} \cdescription

\end{titlepage}
\setcounter{page}{2} % count the title page as page 1


  \section{Learning Objectives}
  \label{objectives}

  \edobject

\section{Required Texts \& Materials}
\label{texts}

The following texts are required. Students are strongly encouraged to
purchase their own copies. Digital copies are available within
\href{https://www.logos.com/product/177582/learning-biblical-hebrew-bundle}{Logos}.
A hard copy of the workbook at least is preferable, since it is filled
with worksheets on which you will want to write.

\begingroup
\renewcommand{\section}[2]{}% temporarily remove the section heading
\begin{thebibliography}{Workbook}% use the longest item in the bibliography

    \bibitem[Grammar]{lbh}
    Karl V. Kutz and Rebekah Josberger.
    \emph{Learning Biblical Hebrew: Reading for Comprehension: An Introductory Grammar}.
    Bellingham, WA: Lexham, 2018. ISBN 978-1683590842.
    Note that LBH has a \href{http://www.learningbiblicalhebrew.com/}{companion site with resources for teachers and students}.

    \bibitem[Workbook]{work}
    Karl V. Kutz and Rebekah Josberger.
    \emph{Learning Biblical Hebrew: A Graded Reader with Exercises}.
    Bellingham, WA: Lexham, 2019. ISBN 978-1683592440.

\end{thebibliography}
\endgroup

\section{Supplementary Texts}
\label{supplementary}

The following reference works are worth owning and consulting.
\cite{rlgs} in particular contains sound advice on core skills like
reading religious texts, writing essays and reviews, revising essays,
making oral presentations, and learning languages.

\begingroup
\renewcommand{\section}[2]{}% temporarily remove the section heading
\begin{thebibliography}{Making Sense}% use the longest item in the bibliography

    \bibitem[BHS]{bhs}
    \emph{Biblia Hebraica Stuttgartensia}.
    Stuttgart: Deutsche Bibel\-ge\-sell\-schaft, 1997. ISBN 978-1598561630.
    This remains the standard critical edition, though it is gradually
    being replaced by the BHQ. A paperback edition, marketed to
    students, is also available; the hardback is much more durable.

	\bibitem[Making Sense]{rlgs} Northey, Margot, Bradford A. Anderson, and Joel N. Lohr.
	\emph{Making Sense in Religious Studies: A Student's Guide to Research and Writing}.
	3rd ed. Toronto: Oxford University Press, 2019. ISBN 978-0199026838.

	\bibitem[SBL2]{sbl2} Collins, Billie Jean, et al.
	\emph{The SBL Handbook of Style}.
	2nd ed. Atlanta: SBL Press, 2014. ISBN 978-1589839649. See the free
	\href{https://www.sbl-site.org/assets/pdfs/pubs/SBLHSsupp2015-02.pdf}{Student Supplement}.

\end{thebibliography}
\endgroup

\section{Course Outline}
\label{outline}

We will adhere to the schedule in \autoref{schedule} as closely as
possible, though the professor reserves the right to adjust it to suit
the needs of the class.

\setcolumncount{4}% set up \sessioncount, \unit{}, \noclass{}, and \reminder{memo}{date} macros
\begin{table}[htbp]% set to `p' to put the schedule on its own page
  \centering%\addfontfeature{ItalicFont=SBLBibLit}
  \begin{tabular}{@{}r<{.}@{ }llr}
    \toprule
    \bfseries Ch & \textbf{Learning Biblical Hebrew: Grammar} & \textbf{Workbook} & \textbf{Date} \\
    \midrule
    1  & The Hebrew Alphabet                                          & E1 (a–b) & 13 June \\
    2  & Hebrew Vowels                                                & E2       & 14 June \\
    3  & Syllables and Reading Hebrew                                 & E3       & 15 June \\
    4  & Gender and Number, Definite Article, Conjunction             & E4, R4   & 16 June \\
    5  & Vowel Changes in Hebrew Nouns                                & E5, R4   & 17 June \\
    6  & Noun and Adjective Function                                  & R6       & 20 June \\
    7  & Constructs, Directional Ending, Prepositions, Interrogatives & E7, R7   & 21 June \\
    8  & Pronominal Suffixes and Review of Definiteness               & E8, R8   & 22 June \\
    9  & Learning to Read Intuitively: Common Patterns in Nouns       & R9–10    & 23 June \\
    10 & Numbers                                                      & R9–10    & 24 June \\
    \bottomrule
  \end{tabular}
  \caption{Schedule of Lessons by Chapter, Exercise (E) \& Reading (R)}
  \label{schedule}
\end{table}

See the AST website for a list of other \href{http://www.astheology.ns.ca/students/academic-dates.html}{important dates}.

\section{Evaluation}
\label{evaluation}

The grade structure for \ccode has the following elements.

\begin{enumerate}
 \item
    Students are expected to read and study all \textbf{lessons} in
    Chapters 1–10 of the \cite{lbh}, and to complete the accompanying
    \textbf{exercises} and Hebrew \textbf{readings} in the \cite{work}.
    Students should use the answer key to correct and evaluate their own
    work. Assignments will be assessed at the end of each week for
    completeness only. Students who are taking the course for credit
    should be prepared to show evidence of their efforts during office
    hours on Fridays, and, if requested, to upload scans (PDFs) or
    photos of any relevant pages \cite[3–34, 79–101]{work}.
 \item
    \textbf{Quizzes} are based on the \cite{lbh} and \cite{work}
    assignments. Between four and six quizzes will be set throughout the
    two-week course. All quizzes are closed book. Students must take
    opportunities accurately to assess what they have learned and
    retained. Continuing education students and auditors are encouraged
    to participate in these exercises fully by scoring and correcting
    their own work.

%%%%%%% For online delivery
% \item
%    Students who are taking the course for credit must use a
%    \textbf{proctor} to administer and attest their quizzes. Select
%    someone over the age of 18 to act as your proctor. Instruct them to
%    email your instructor
%    (\href{mailto:ddriver@astheology.ns.ca}{ddriver@astheology.ns.ca})
%    to receive electronic copies of all six quizzes for this course.
%    When you are ready to take each quiz, your proctor will provide you
%    with a printed copy of the test materials. Be sure that you include
%    compensation for them to do so, as appropriate. When you are
%    finished with each quiz, the proctor will verify that it was taken
%    as prescribed, attest it by signing their name under your name, and
%    make sure that you submit it online. Please upload each quiz to the
%    correct assignment slot on Teams. You will obviously need to scan
%    each one as a PDF (preferred) or take a photo of each page with your
%    cell phone.
\end{enumerate}

The breakdown for the semester's total work is shown in
\autoref{distribution}.

\begin{table}[htbp]
  \centering
  {\lining
  \begin{tabular}{lr}
    \toprule
    Workbook & 50\% \\
    Quizzes  & 50\% \\
    \bottomrule
  \end{tabular}}
  \caption{Distribution of Grades}
  \label{distribution}
\end{table}

\ProvidesFile{grades.tex}[2023/09/06 v2.5 -- Course policy]

AST's \href{https://www.astheology.ns.ca/students/resources.html}{Academic
Calendar} provides guidelines and criteria for academic assessment. Marks are
assigned by letter grade using these benchmarks. Note that graduate and
undergraduate programs are currently assessed on a different scales at AST (see
\autoref{grade-syst}). For graduate students, a grade of C is deemed
``Unsatisfactory,'' whereas for undergraduate (BTh) students, a C is
``Satisfactory'' and the grade of D is possible for ``Marginal'' performance.
Other grades have different definitions, too.

\begin{table}[htbp]
 \centering
 {\lining\footnotesize
 \begin{tabular}{lll|ll}
   \toprule
               & \multicolumn{2}{c|}{Graduate Courses}   & \multicolumn{2}{c}{Undergraduate Courses} \\
               \cline{2-5}
   A+          & 94--100    & Exceptional               & 90--100    & Excellent                 \\
   A           & 87--93     & Outstanding               & 85--89     & Excellent                 \\
   A\char"2212 & 80--86     & Excellent                 & 80--84     & Excellent                 \\ [1ex]
   B+          & 77--79     & Good                      & 75--79     & Good                      \\
   B           & 73--76     & Acceptable                & 70--74     & Good                      \\
   B\char"2212 & 70--72     & Marginally Acceptable     &            &                           \\ [1ex]
   C           & 60--69     & Unsatisfactory            & 60--69     & Satisfactory              \\
   D           &            &                           & 50--59     & Marginal                  \\
   F           & 0--59      & Failure                   & 0--49      & Failure                   \\
   FP          & 0          & Failure due to Plagiarism & 0          & Failure due to Plagiarism \\
   \bottomrule
 \end{tabular}}
 \caption{Outline of Grades at AST}
 \label{grade-syst}
\end{table}

\ProvidesFile{other.tex}[2018/0/03 v2.4 -- Course policy]

\section{Other Course Policy}
\label{policy}

Late work will not be accepted, except in genuinely extenuating
circumstances. Students must submit something before the deadline if
they wish to receive credit. Unless I state otherwise, written
assignments are to be uploaded by 11:59 \PM\ on the date indicated.

Essay submissions must be typewritten, double-spaced, and formatted as
PDFs. They should be free from error. In this course they should follow
either \emph{The SBL Handbook of Style} (refer to the online
\href{https://www.sbl-site.org/assets/pdfs/pubs/SBLHSsupp2015-02.pdf}{
Student Supplement}), or \emph{The Chicago Manual of Style}, per AST's
policy for major papers. As a reminder, AST also upholds an Inclusive
Language Policy.

\href{http://www.eerdmans.com/Pages/Item/59043/Commentary-Statement.aspx}{Plagiarism},
if \href{https://www.theguardian.com/world/2013/feb/09/german-education-minister-quits-phd-plagiarism}{detected},
will result in failure of the course.

Students should request permission to record a class or lecture. If
permission is granted, or if recordings are provided (as in the case of
an online course), I stipulate that all recordings be for personal use
only. They may not be shared or distributed.

If you have abilities or disabilities that require modifications to the
assessment process or other aspects of this course, please advise the
course instructor as soon as possible.

Finally, I encourage the conscientious use of laptops, tablets, and
other technology in my classes. In classroom settings, realize that, as
\href{http://dx.doi.org/10.1016/j.compedu.2012.10.003}{cognitive
psychologists have demonstrated}, ``laptop multitasking hinders
classroom learning for both users and nearby peers.'' Do your part to
foster an environment of open dialogue by honouring the presence of your
classmates. In online settings, consider both the physical environment
in which you choose to work and the virtual environment that you help
create through your participation in various forums. Let your engagement
in this course be marked by rigour and charity alike.


\section{Bibliography}
\label{bib}

Among the many Hebrew grammars and resources, note the following
especially:

\begin{itemize}
  \item Guides \& Aids

    \begin{itemize}
      \item
        Van Pelt, Miles. \emph{English Grammar to Ace Biblical Hebrew}.
        Grand Rapids: Zondervan, 2010. At AST: PE 1130 H5 V35 2010
      \item
        Einspahr, Bruce. \emph{Index to Brown, Driver \& Briggs Hebrew Lexicon}.
        Chicago: Moody Press, 1976. At AST: PJ 4833 B683 E35
      \item
        Landes, George M. \emph{Building Your Biblical Hebrew
        Vocabulary: Learning Words by Frequency and Cognate}.
        Atlanta: SBL, 2001. At AST: PJ 4845 L25 2001
      \item
        Scott, W.\,R., and H.\,P. Rüger. \emph{A Simplified Guide to BHS:
        Critical Apparatus, Masora, Accents, Unusual Letters \& Other
        Markings}. 3rd edition. N. Richland Hills, TX: BIBAL, 2007.
        At AST: BS 715 1977f
      \item
        Kelley, Page H., Daniel S. Mynatt, and Timothy G. Crawford.
        \emph{The Masorah of Biblia Hebraica Stuttgartensia:
        Introduction and Annotated Glossary}.
        Grand Rapids: Eerdmans, 1998. At AST: BS 718 K38 1998
    \end{itemize}

  \item Introductory Grammars

    \begin{itemize}
      \item
        Cook, John A. and Robert D. Holmstedt.
        \emph{Beginning Biblical Hebrew: A Grammar and Illustrated Reader}.
        Grand Rapids: Baker Academic, 2013.
        Study aids are available at \href{http://www.bakerpublishinggroup.com/books/beginning-biblical-hebrew/5629/students/esources}{bakeracademic.com}.
        At AST: PJ 4567.3 C66 2013
      \item
        Kittel, Bonnie, Vicki Hoffer, and Rebecca Abts Wright.
        \emph{Biblical Hebrew: Text and Workbook} [and Audio CD].
        2nd edition. New Haven: Yale, 2004.
      \item
        Martin, J.\,D. \emph{Davidson’s Introductory Hebrew Grammar}.
        27th edition. Edinburgh: T\&T Clark, 1993. At AST: PJ 4567 D37 1962
      \item
        Muraoka, Takamitsu.
        \emph{A Biblical Hebrew Reader: With an Outline Grammar}.
        Leuven: Peeters, 2017. At AST: PJ 4567.3 M87 2017
      \item
        Pratico, Gary and Miles Van Pelt. \emph{Basics of Biblical Hebrew Grammar}.
        Grand Rapids: Zondervan, 2001. At AST: PJ 4567.3 P73 2007
      \item
        Seow, C.\,L. \emph{A Grammar for Biblical Hebrew}.
        Nashville: Abingdon, 1995. At AST: PJ 4567 S424 1995
      \item
        Weingreen, J. \emph{A Practical Grammar for Classical Hebrew}.
        Oxford: Clarendon, 1959. At AST: PJ 4567 W4 1959
    \end{itemize}


  \item Syntaxes \& Reference Grammars

    \begin{itemize}
      \item
      	Arnold, Bill T., and John H. Choi. \emph{A Guide to Biblical Hebrew Syntax}.
      	Cambridge: Cambridge University Press, 2003. At AST: PJ 4701 A76 2003
      \item
      	Waltke, Bruce K., and M. O’Connor. \emph{An Introduction to Biblical Hebrew Syntax}.
      	Winona Lake, IN: Eisenbrauns, 1990.
      \item
        Williams, Ronald J., and  Beckman, John C. \emph{Williams' Hebrew Syntax}.
        Toronto: University of Toronto Press, 2007. At AST: PJ 4701 W5 2007
      \item
        Gesenius, Wilhelm, E. Kautsch and A.\,E. Cowley. \emph{Hebrew Grammar}.
        Oxford: Clarendon, 1910. \textbf{GKC} remains a standard Hebrew
        reference grammar in English, even though Cowley's translation
        -- of the 28th german edition -- is now over 100 years old. At
        AST: PJ 4564 G5 1910
      \item
      	Joüon, P., and T. Muraoka. \emph{A Grammar of Biblical Hebrew}.
      	2nd edition. Subsidia Biblica 27. Rome: Pontifical Biblical
      	Institute, 2008. One of the most complete and up-to-date
      	Hebrew grammars in English, \textbf{JM} (1st ed. 1991) was
      	revised from a French work first published by Paul Joüon in
      	1923. At AST: PJ 4567 J7613 1993
    \end{itemize}

  \item Lexicons

    \begin{itemize}
      \item
        Brown, Francis, S.\,R. Driver and Charles A. Briggs. \emph{The
        Brown-Driver-Briggs Hebrew and English Lexicon}. Peabody, Mass:
        Hendrickson, 2004. The \textbf{BDB} was originally published in 1906.
        On \href{https://archive.org/details/hebrewenglishlex00geseuoft}{archive.org}
        and at AST: PJ 4833 B68 1996
      \item
        Köhler, Ludwig, and Walter Baumgartner. \emph{The Hebrew
        and Aramaic Lexicon of the Old Testament}. Leiden: Brill,
        1994--2000.
        At AST: PJ 4833 K61813 1994
      \item
        Holladay, William L. \emph{A Concise Hebrew and Aramaic Lexicon
        of the Old Testament, Based Upon the Lexical Work of Ludwig
        Koehler and Walter Baumgartner}.
        Grand Rapids: Eerdmans, 1971.
        At AST: PJ 4833 H6 1971
      \item
        Clines, David J.\,A. \emph{The Dictionary of Classical Hebrew}.
        Sheffield: Sheffield Academic Press, 1993--2011.
        At AST: PJ 4833 D53 1993
    \end{itemize}

\end{itemize}


\end{document}
